\section{Purpose and Scenario}

\begin{memobox}
\textbf{Memorandum --- Commandant Isabella Ricci, Polizia Postale}\\
\textbf{To:} Digital Forensics Unit --- All Roles\\
\textbf{Subject:} New Engagement --- NovaDynamics S.p.A., Operation SILENT BUCKET

\vspace{0.3cm}
Team; three months after the ransomware incident, NovaDynamics S.p.A.\ completed a
migration of its production workloads to a cloud provider, seeking better resilience.
Yesterday, their CISO Dr.\ Silvia Valenti reported that a batch of confidential client
contracts has surfaced for sale on an underground forum. The affected data was stored in
a cloud object storage bucket used by NovaDynamics' billing service (IaaS-hosted VM) and
a separate customer-facing PaaS application.

\vspace{0.2cm}
NovaDynamics has authorized limited API-level access to their cloud environment, mediated
through their cloud service provider, under a new Forensic Analysis Assignment dated 18
May 2026. As in past engagements, physical access to any hardware is out of question.

\vspace{0.2cm}
I expect the team to determine what happened, when, and whether the evidence available to
us is sufficient to support a judicial case, given the peculiar constraints of this
environment. Coordinate explicitly, and  document exactly
\emph{what you could and could not verify}.

\vspace{0.2cm}
\hfill Commandant Isabella Ricci --- Polizia Postale, Digital Forensics Unit
\end{memobox}

\noindent
This laboratory scope covers: the IaaS/PaaS/SaaS distinction and its impact on evidence availability; interaction with a cloud object store via API; interpretation of provider-supplied event logs and their inherent gaps; identification of the authoritative
(``primary'') copy among redundant/replicated objects; and the legal and trust implications of relying on cloud-provider-mediated evidence.

\subsection{Investigation Objectives (from the Engagement Letter)}
NovaDynamics S.p.A.\ formally requests that the team:
\begin{itemize}[leftmargin=1.2cm]
	\item Determine which cloud service model (IaaS/PaaS/SaaS) applies to each affected
	      component, and what evidence is realistically obtainable from each.
	\item Reconstruct a timeline of the exfiltration event using the log material the cloud
	      provider has made available.
	\item Identify which of the several copies/versions of the affected storage object is
	      authoritative, and document this in the chain of custody.
	\item Assess the legal admissibility of cloud-provider-mediated evidence, and discuss the
	      jurisdictional complications raised by the case.
	\item Deliver a short technical report addressing the above, including an explicit
	      statement of the investigation's limitations.
\end{itemize}

\subsection{What We Already Have, and What We Still Need to Find Out}

In a disk or memory investigation, the evidence is a single artefact (a disk image, a memory dump) that you acquire once and then examine offline, at your own pace, with full control over the copy. In a cloud investigation, there is no single artefact to acquire: what NovaDynamics' cloud provider gives you is a set of \emph{exports}, such as logs, metadata listings, billing reports, each produced by a different internal system of the provider, on a different schedule, with
different blind spots.\\
Your job is not just to \emph{read} these exports, but to figure out \emph{what they can and cannot tell you}, and to fill the gaps with reasoning, not with more raw
disk bytes (because there are none you can access).
\medskip

\noindent
That is exactly why this laboratory has two distinct kinds of activity:

\begin{itemize}[leftmargin=1.2cm]
	\item \textbf{Sections 3--4 (hands-on, ``live'')}: you will interact directly with a
	      \emph{local, harmless stand-in} for the real bucket, using the same commands and the
	      same protocol (S3 API) that a real investigator would use against the real cloud
	      provider.\\
		  The goal here is not to find evidence but to build an intuition for \emph{what kind of information the API itself can give you} (object listings, metadata, hashes), so that later, when reading the real exports, you understand exactly how that information was originally produced and what its natural limitations are.
	\item \textbf{Sections 5--6 (analysis, ``the actual case'')}: you will examine the exports
	      below, which represent what NovaDynamics' cloud provider actually handed over for this
	      case. This is where the real investigation happens.
\end{itemize}

\subsection{Evidence \& Access Manifest}

\begin{table}[H]
	\centering
	\begin{tabularx}{\textwidth}{l | l | X}
		\toprule
		\textbf{Artefact} & \textbf{Format} & \textbf{Description} \\
		\midrule
		\texttt{cloudtrail\_events.json}    & JSON                & Management/data-plane event log export, provided by the CSP \\
		\texttt{cloudtrail\_digest.json}    & JSON                & Digest file accompanying the log export, used by the CSP to attest log integrity \\
		\texttt{s3\_object\_versions.json}  & JSON                & Metadata export: versions, timestamps, region tags for the affected object \\
		\texttt{data\_transfer\_metrics.json} & JSON              & Independent billing/data-transfer-out report, hourly granularity \\
		Live sandbox access                 & MinIO (Docker)      & A local S3-compatible emulation of the bucket, for direct hands-on interaction \\
		\bottomrule
	\end{tabularx}
	\caption{Evidence and access manifest for Operation SILENT BUCKET.}
\end{table}

\noindent
The second artefact, \texttt{cloudtrail\_digest.json}, deserves a note here because it is the mechanism a real cloud provider uses to attest that a log export has not been altered after it was produced. Real AWS CloudTrail digest files chain each digest to the hash of the previous one and are signed with the provider's private key, so that a recipient can verify both the content and the origin of the logs. The file supplied for this laboratory is a simplified version of that format; you will examine what it does and does not let you prove in Section 8.
\medskip

\begin{warningbox}
All JSON artefacts must be treated as evidence exports. Verify the integrity of each file
before analysis by computing its SHA-256 hash and recording it in your chain-of-custody log.\\
No artefact may be modified in place; work only on copies.\\
Note that a hashing/query tool such as \texttt{jq} (Section 5) does not, on its own, make an
analysis forensically sound. Keep this
distinction in mind: it will resurface in Section 8.
\end{warningbox}
